\section{C++ survival components: preserving the event, clock and output}
\label{sec:exec24-longitudinal-cxx}

C++ survival components answer several different questions inside a banking
workflow. A collections team may want the observed probability of default by
month twelve, the default probability in a hypothetical world without prepayment,
or the effect of a sequence of contact decisions. These are different targets.
A fast predictor of recorded survival time does not choose among them. The
proposed product must preserve the event dictionary, observation and treatment
clocks, competing transitions, selection law and uncertainty target before it
chooses a computational engine. This comparison examines four pinned lineages:
XGBoost, CatBoost, ranger and aorsf. The evidence consists of authenticated primary
source and documentation, not execution of these packages or their tests.
Discovery remains bounded; neither four candidates nor a repository search
establishes exhaustive C++ coverage.~\citep{exec24_longcxx_xgboost,exec24_longcxx_catboost,exec24_longcxx_ranger,exec24_longcxx_aorsf}

\subsection{The estimand precedes the software label}

Let $T$ be the first transition time and $J$ its type, with $J=D$ for default
and $J=P$ for prepayment. The default cumulative incidence is
$F_D(t)=\Pr(T\leq t,J=D)$. It is not generally the complement of survival
obtained by treating prepayments as ordinary censoring. In an authored finite
cohort of ten accounts, four prepay at time one, two default at time two and
four remain open. With these complete observations, the default incidence is
$2/10$. The product of surviving all earlier transitions and the conditional
default fraction is $(6/10)(2/6)=1/5$. Treating the four prepayments as censored
instead gives the default-only complement $2/6=1/3$. This is a different
functional, not an arithmetic refinement of the same answer. At a stipulated
USD~1,000 default loss, the two expected-loss calculations are USD~200 and
USD~$1{,}000/3$. Neither a different language nor more trees resolves the
target substitution. A hypothetical intervention eliminating prepayment also
needs its own scientific assumptions; censoring a record is not such an
intervention.~\citep{exec24_longcxx_reference}

Observation intervals need equally careful interpretation. If account $i$
enters observation at $L_i$ and exits at $U_i$, the risk indicator for a
counting-process analysis is $Y_i(t)=1\{L_i<t\leq U_i\}$. For the three
authored intervals $(0,1]$, $(1,2]$ and $(0,2]$, with a default on the first
account at time one, the correct risk set contains two accounts. Its unit-risk
hazard increment is $1/2$, compared with $1/3$ if every terminal-time row is
treated as present from zero. Bounds on an unknown failure time are different
from entry and exit times: an interval-censored observation says that an event
occurred within a range, whereas delayed entry says when an account becomes
observable. An adapter must refuse their silent interchange. Repeated rows also
do not by themselves provide sequential exchangeability, positivity or
cluster-valid uncertainty for a treatment policy.~\citep{exec24_longcxx_reference}

\subsection{Boosted Cox and accelerated failure time objectives}

The selected XGBoost commit is
\texttt{c914bd00a946081c96e2f7aaa0e298f38b7b17e1}, dated 7 October 2026.
Its root build declares version 3.5.0; this identifies inspected development
bytes, not a supported binary release. The Cox objective sorts absolute labels,
uses positive labels for observed events and negative labels for censoring,
and exponentiates the margin for prediction. That exponentiated output is a
relative risk score. It does not contain a baseline survival curve or a monetary
default probability. The inspected routine multiplies each final gradient and
Hessian by its row weight, while its risk-set exponential sum is unweighted.
This source contract must not be relabelled as an inverse-probability weighted
longitudinal estimator without a separate derivation.~\citep{exec24_longcxx_xgboost}

XGBoost's AFT documentation instead specifies
$\log T=m(x)+\sigma Z$, with normal, logistic or extreme-value choices for
$Z$. Its ranged labels encode exact, right-, left- or interval-censored failure
times. The examples use zero as a left bound and positive infinity as the right
bound for right censoring. The inspected CPU and CUDA kernels call the same
distribution-specific derivative helpers and apply row weights. The prediction
transform returns $\exp(m(x))$, while the evaluation transform retains the raw
margin for the AFT loss. Declared CUDA kernels are evidence of implementation
text, not observed GPU conformance, speed or numerical parity. The pinned
tutorial routes ranged labels through the native training and data-matrix
interfaces and states a limitation for its scikit-learn interface. These
statements are tied to the acquired version.~\citep{exec24_longcxx_xgboost}

The AFT output requires its distribution and scale to be useful. Consider an
authored logistic AFT law with $\exp(m)=2$ and horizon eight. For scale one,
the failure probability is $F(\log 4)=4/5$; for scale two it is
$F(\log 2)=2/3$, where $F(z)=e^z/(1+e^z)$. Both models return the same
exponentiated location. A consumer that interprets that location as a probability
cannot distinguish them. The proposed result therefore carries the raw margin,
transformed location, distribution, scale, time unit and requested horizon as
separate fields.~\citep{exec24_longcxx_reference}

The CatBoost commit
\texttt{1ac99c589b54e91a349dd04a049f8ec1e1572326}, also dated 7 October
2026, declares Python package version 1.2.10. Its documentation presents Cox
and SurvivalAft objectives and marks their GPU support unavailable in the
inspected table. The displayed Cox expression uses $a_i$ inside an exponential,
and the AFT helper uses $(\log t-a)/\sigma$. These mathematical roles require
more precise output typing than the documentation's common phrase \emph{hazard
rates}. The raw objective margin is not an absolute hazard or probability.
CatBoost's inspected AFT routine uses $-1$ as an unbounded endpoint, with a
separate lower-end branch for left censoring. XGBoost's documented zero and
infinity conventions therefore cannot simply be copied into a CatBoost target
array. End-to-end target validation and the public prediction transformation
still need a frozen wrapper review.~\citep{exec24_longcxx_catboost}

Both projects clip or regularize AFT derivatives and provide limiting branches
for numerically extreme inputs. These choices are part of the computation,
not proofs of statistical calibration. CatBoost's selected Cox and AFT derivative
signatures leave their weight argument unused, but the inspected Cox caller
also passes weights into subsequent leaf accumulation. The appropriate finding
is that weighted-objective equivalence is unqualified: a local unused argument
does not establish that the entire training system ignores weights. A prototype
must compare the complete selected fitting path with an independently specified
weighted target.~\citep{exec24_longcxx_xgboost,exec24_longcxx_catboost}

\subsection{A tie fixture that can falsify silent equivalence}

Source review supplies a concrete acceptance question before a costly benchmark.
Take three unit risk scores, with two observed events tied at time one and a
third account censored at time two. For margins $f_i$, the Breslow negative
log-likelihood contribution is
\[
 \ell(f)=-f_1-f_2+2\log\{e^{f_1}+e^{f_2}+e^{f_3}\}.
\]
At zero margins its gradient is $(-1/3,-1/3,2/3)$. Its components sum to zero,
as required by invariance to a common additive margin. An independent rational
trace of the inspected XGBoost loop, in event/event/censor order, gives
$(-2/3,-1/3,2/3)$: the first tied event is processed before the second event's
reciprocal-risk contribution enters the running accumulator. The sum is
$-1/3$. The analogous inspected CatBoost loop removes the previous exponential
at each positive event and gives negative-log-likelihood-sign scores
$(-2/3,-1/6,5/6)$. Their sum is zero, but the two tied events are not symmetric
and the vector differs from the stated Breslow derivative. These are abstract
source-loop traces, not native benchmark results or a claim that every release,
wrapper or training regime behaves identically.~\citep{exec24_longcxx_xgboost,exec24_longcxx_catboost,exec24_longcxx_reference}

The next justified experiment is narrow: freeze the complete dispatch path,
evaluate equal-time permutations and event/censor ties, compare finite-difference
derivatives with an independent likelihood, and retain both the returned
gradients and row identities. Non-unit weights and extreme margins form separate
fixtures. A failing fixture would justify a restricted regime, another engine,
or a reviewed repair; it would not justify silently changing the scientific
likelihood. Until that experiment is admitted and performed, these findings
remain source-level concerns and acceptance criteria.

\subsection{Survival forests and transformations that change the decision}

ranger is pinned to
\texttt{dafa5db57972dd49c808db881e6e6cd5c4ebdc4a}, dated 16 January
2026; its R description declares 0.18.0. The inspected C++ leaf routine sums
$d_j/n_j$ to construct a cumulative hazard. The forest averages those hazards,
and the inspected R prediction wrapper returns survival as the exponential of
their negative. Its OOB error uses one minus concordance on summed cumulative
hazards for accounts with at least one OOB tree. That is a ranking diagnostic,
not an interval for a causal effect or a certification of horizon calibration.
The wrapper rejects nonbinary event coding in the selected survival path, and
its two standard-error choices explicitly reject survival forests. Case weights
are documented as sampling weights for bootstrap or subsampling selection.
These contracts are more specific than the labels \emph{weighted forest} or
\emph{uncertainty available}.~\citep{exec24_longcxx_ranger}

aorsf is pinned to
\texttt{085db7062d995eeafa4e210481d66af2eb02b8dc}, dated 5 February
2026; its description declares 0.1.6. Its documented accelerated survival
control uses a one-iteration Cox fit to construct split directions. The C++
delegate distinguishes that path from further iterations. This is a tree-building
choice, not a claim that the final forest is a converged Cox regression.
The selected leaf code maintains a product-limit survival estimate and an
additive cumulative-hazard estimate separately, with weighted event and risk
counts. The wrapper documents normalization of supplied weights to the number
of observations. Risk, survival, cumulative hazard, mortality and time are
distinct prediction types. Predictions beyond maximum follow-up are documented
as constant if boundary checks are disabled; a flat continuation is not evidence
about an unobserved future.~\citep{exec24_longcxx_aorsf}

For an authored leaf with one event among two at risk, the product-limit
survival is $1/2$, whereas the cumulative hazard is $1/2$ and its exponential
survival is $e^{-1/2}\simeq0.606531$. The corresponding risks are $0.5$ and
approximately $0.393469$. With a stipulated USD~450 premium and USD~1,000 loss,
the margins are USD~$-50$ and approximately USD~56.53. These are deliberately
simple constructed outputs; they are not observed errors or bank estimates.
They show why a consumer must retain the estimator and transformation, even
when both fields are called survival. Aggregation creates another distinction:
for two authored hazards $(0,\log4)$, exponentiating the negative average gives
$1/2$, while averaging their exponentiated negatives gives $5/8$. An adapter
must reproduce the declared aggregation order rather than substitute a familiar
formula.~\citep{exec24_longcxx_reference}

\subsection{Integration, adoption and justified rejection}

The proposed first route retains original engines behind typed workers. XGBoost
and CatBoost are conditional boosted-prediction alternatives once their objective,
endpoint, margin, weight, tie and environment contracts pass selected fixtures.
ranger's standalone C++ library or CLI and its R wrapper are separate packaging
choices; the inspected build declares C++14 and library/CLI targets. Its primary
headers explicitly distinguish the MIT C++ core from the GPL-3 R package.
XGBoost and CatBoost have Apache-2.0 root declarations. aorsf declares MIT and
R/Rcpp/RcppArmadillo dependencies. These direct declarations do not clear every
dependency, model artifact or redistribution. Rights review must bind the actual
chosen distribution.~\citep{exec24_longcxx_xgboost,exec24_longcxx_catboost,exec24_longcxx_ranger,exec24_longcxx_aorsf}

For aorsf, an original R worker preserves the formula, validation and prediction
semantics while using its C++ core. Extracting that core into a new standalone
port would require recreating the R-side contract and numeric dependencies;
the proposal has not demonstrated a benefit sufficient to justify that work.
For ranger, direct C++ embedding remains an alternative to an R worker when
deployment requirements and rights justify the separate interface. Whole-library
rewrites are rejected for the present proposal because they add parity and
maintenance obligations without demonstrated need. Native workers, a dedicated
longitudinal R/Python estimator and an independent finite-law reference remain
substantive alternatives, selected according to the target rather than language
preference.~\citep{exec24_longcxx_aorsf,exec24_longcxx_ranger}

The CLI and MCP proposal should expose a versioned request containing population,
time zero, availability clock, event dictionary, competing-event interpretation,
entry/censoring representation, treatment history, weights and their meaning,
sampling units, horizons, options and artifact hashes. The result must preserve
accepted and rejected row identities, objective and output type, baseline or
distribution parameters, convergence/numerical warnings, uncertainty method,
timing and complete request costs. Refusal is a typed result. A finite output,
successful process exit or replayed artifact must not be promoted to causal
identification or qualified uncertainty. These are proposed interfaces; no new
native adapter or sibling integration is implemented by this review.

The method owner must first decide whether a public or generated study can
identify the proposed target. The adapter owner then freezes version, wrapper
and conversion fixtures; the rights owner resolves the actual distribution;
the evaluation owner preregisters held-out targets, ties, censoring and competing
events, row permutations, dependent units, failure denominators and full costs.
No purchase, vendor trial or expensive bank experiment is necessary to write
those contracts. Candidate aliases, unreturned discovery tails, uninspected
delegates, cluster-calibrated uncertainty and native performance remain recorded
follow-ups. This bounded comparison improves the proposal's executable evidence
and actions without declaring C++ longitudinal coverage complete.
